\chapter{Resultados Experimentales}
\label{chap:ejecucion-obtencion-resultados}


\section{Configuración Final del AG}

A partir del proceso de calibración presentado en las subsecciones anteriores ---donde se evaluaron de forma independiente el operador de cruce, el número de generaciones y la probabilidad de mutación--- se obtuvo una configuración final del algoritmo genético que equilibra de manera adecuada la calidad de las soluciones y el costo computacional. Esta configuración fue utilizada en todas las comparaciones realizadas posteriormente contra el método SPT en el Capítulo~7.

En síntesis, la calibración determinó que:

\begin{itemize}
    \item El operador de cruce \textit{single\_point} ofrece el mejor desempeño en un horizonte amplio de 480 minutos, superando consistentemente a \textit{two\_points}.
    \item Un número de \textbf{300 generaciones} es suficiente para que el algoritmo converja hacia soluciones de alta calidad, mientras que valores mayores no aportan mejoras significativas y pueden incluso degradar el rendimiento.
    \item Una probabilidad de mutación del \textbf{15\%} logra el mejor equilibrio entre exploración y explotación del espacio de búsqueda, evitando tanto la convergencia prematura (como sucede con 5\%) como la pérdida de estabilidad evolutiva (observada con 25\%).
\end{itemize}

Con base en estos resultados, la configuración final del algoritmo genético utilizada en las simulaciones comparativas es la siguiente:

\begin{table}[htbp]
\centering
\small
\caption{Configuración final seleccionada para el algoritmo genético.}
\label{tab:config-final-ag}
\begin{tabular}{@{}lc@{}}
\toprule
\textbf{Parámetro} & \textbf{Valor seleccionado} \\
\midrule
Operador de cruce & \textit{single\_point} \\
Número de generaciones & 300 \\
Tamaño de población & 100 \\
Método de selección & \textit{rank} \\
Padres por generación & 10 \\
Padres preservados & 5 \\
Probabilidad de mutación & 15\% \\
Semilla aleatoria & 42 \\
\bottomrule
\end{tabular}
\end{table}

Esta configuración proporciona un desempeño robusto en términos de tareas ejecutadas, carga total y estabilidad del proceso evolutivo. Además, se mantiene dentro de límites razonables de tiempo de ejecución, lo que facilita su aplicación práctica en un taller industrial. Los experimentos del capítulo siguiente emplean exclusivamente esta configuración, garantizando que la comparación frente al método SPT se realice bajo condiciones justas y reproducibles.
